% Generated from lessons/0007-4-geometric-applications-definite-integrals.html; edit the HTML source, then rerun the converter.
\section{定积分在几何计算中的应用}
\lessonmark{定积分在几何计算中的应用}

这一节的核心不是背很多公式，而是学会把几何量切成“微小片段”，再用定积分累加。面积、弧长、体积、旋转曲面面积、曲率都遵循同一个动作：选变量、写微元、定上下限。

\subsection*{本节做题流程}

\begin{conceptbox}
\textbf{先画图}
确定谁在上、谁在下，谁是外半径、谁是内半径，是否要分段。
\end{conceptbox}

\begin{conceptbox}
\textbf{选微元}
面积用小矩形，体积用薄片或圆柱壳，弧长用微小线段。
\end{conceptbox}

\begin{conceptbox}
\textbf{检查符号}
面积、长度、体积非负；若公式可能带方向，必须取绝对值或分段。
\end{conceptbox}

\begin{notebox}
这一节最容易错在“范围”和“分段”：交点、参数区间、极坐标角度范围、旋转半径都要先确认。
\end{notebox}

\subsection*{一、平面图形面积}

\subsubsection*{直角坐标}

\[
      S=\int_a^b |f(x)|\,dx,\qquad
      S=\int_a^b |f(x)-g(x)|\,dx.
    \]

若已经知道 \(f(x)\ge g(x)\)，则可直接写

\[
      S=\int_a^b [f(x)-g(x)]\,dx.
    \]

\subsubsection*{参数方程}

若曲线由 \(x=x(t),y=y(t)\) 给出，可用

\[
      S=\int y\,dx=\int_{\alpha}^{\beta}y(t)x'(t)\,dt
    \]

配合方向判断正负。闭合曲线通常要按图形方向取绝对值或利用对称性。

\subsubsection*{极坐标}

\[
      S=\frac12\int_{\alpha}^{\beta}r^2(\theta)\,d\theta.
    \]

这是把扇形面积微元 \(\frac12r^2d\theta\) 累加。

\subsection*{二、弧长}

\begin{center}
\small
\begin{tabularx}{\linewidth}{|>{\raggedright\arraybackslash}X|>{\raggedright\arraybackslash}X|}
\hline
\textbf{曲线形式} & \textbf{弧长公式} \\ \hline
\(y=f(x),\ a\le x\le b\) & \(\displaystyle L=\int_a^b\sqrt{1+[f'(x)]^2}\,dx\) \\ \hline
\(x=x(t),y=y(t),\ \alpha\le t\le\beta\) & \(\displaystyle L=\int_\alpha^\beta\sqrt{[x'(t)]^2+[y'(t)]^2}\,dt\) \\ \hline
\(r=r(\theta),\ \alpha\le\theta\le\beta\) & \(\displaystyle L=\int_\alpha^\beta\sqrt{r^2(\theta)+[r'(\theta)]^2}\,d\theta\) \\ \hline
\end{tabularx}
\end{center}

\textbf{证明思路：}小段曲线近似为直线段，长度微元满足 \(ds=\sqrt{dx^2+dy^2}\)。

\subsection*{三、体积}

\subsubsection*{截面积法}

若垂直于 \(x\) 轴的截面积为 \(A(x)\)，则

\[
      V=\int_a^b A(x)\,dx.
    \]

\subsubsection*{旋转体：圆盘与垫片}

\[
      V=\pi\int_a^b \bigl(R^2(x)-r^2(x)\bigr)\,dx.
    \]

这里 \(R\) 是外半径，\(r\) 是内半径。

\subsubsection*{旋转体：圆柱壳}

\[
      V=2\pi\int_a^b \text{半径}\cdot\text{高度}\,dx.
    \]

绕 \(y\) 轴旋转时，常见形式是 \(V=2\pi\int_a^b x f(x)\,dx\)。

\subsection*{四、旋转曲面面积}

曲线 \(y=f(x)\ge0\) 绕 \(x\) 轴旋转一周，曲面面积为

\[
      S=2\pi\int_a^b y\sqrt{1+(y')^2}\,dx.
    \]

若参数方程 \(x=x(t),y=y(t)\) 绕 \(x\) 轴旋转，则

\[
      S=2\pi\int_\alpha^\beta y(t)\sqrt{[x'(t)]^2+[y'(t)]^2}\,dt.
    \]

绕 \(y\) 轴时，把半径 \(y\) 换成 \(x\)。

\subsection*{五、曲率与曲率半径}

曲率衡量曲线弯曲得有多快。曲率越大，曲线越弯；曲率半径 \(\rho=1/\kappa\)。

\begin{center}
\small
\begin{tabularx}{\linewidth}{|>{\raggedright\arraybackslash}X|>{\raggedright\arraybackslash}X|}
\hline
\textbf{形式} & \textbf{曲率公式} \\ \hline
\(y=y(x)\) & \(\displaystyle \kappa=\frac{|y''|}{(1+(y')^2)^{3/2}}\) \\ \hline
\(x=x(t),y=y(t)\) & \(\displaystyle \kappa=\frac{|x'y''-y'x''|}{([x']^2+[y']^2)^{3/2}}\) \\ \hline
\(r=r(\theta)\) & \(\displaystyle \kappa=\frac{|r^2+2(r')^2-rr''|}{(r^2+(r')^2)^{3/2}}\) \\ \hline
\end{tabularx}
\end{center}

\subsection*{六、课后题训练}

下面按教材第 281-283 页习题 1-19 整理。题面完整保留，提示给出主要切入点。

\begin{exercisebox}
\textbf{习题 1 \badge{平面面积}}

求下列曲线所围的图形面积：

\[
      \begin{aligned}
      &(1)\ y=\frac1x,\ y=x,\ x=2;\\
      &(2)\ y^2=4(x+1),\ y^2=4(1-x);\\
      &(3)\ y=x,\ y=x+\sin^2x,\ x=0,\ x=\pi;\\
      &(4)\ y=e^x,\ y=e^{-x},\ x=1;\\
      &(5)\ y=|\ln x|,\ y=0,\ x=0.1,\ x=10;\\
      &(6)\ x=2t-t^2,\ y=2t^2-t^3,\ 0\le t\le2;\\
      &(7)\ x=a\cos^3t,\ y=a\sin^3t,\ 0\le t\le2\pi;\\
      &(8)\ r=a\theta,\ \theta=0,\theta=2\pi;\\
      &(9)\ r=ae^\theta,\ \theta=0,\theta=2\pi;\\
      &(10)\ r=a\cos\theta+b,\quad b\ge a>0;\\
      &(11)\ r=3\cos\theta,\ r=1+\cos\theta,\quad -\pi/3\le\theta\le\pi/3;\\
      &(12)\ r^2=a^2\cos2\theta;\\
      &(13)\ r=a\cos2\theta;\\
      &(14)\ x^3+y^3=3axy;\\
      &(15)\ x^4+y^4=a^2(x^2+y^2).
      \end{aligned}
      \]

\begin{hintbox}
直角坐标题先找交点；参数和极坐标题优先利用对称性。
\end{hintbox}
\end{exercisebox}

\begin{exercisebox}
\textbf{习题 2 \badge{面积最值}}

求由抛物线 \(y^2=4ax\) 与过其焦点的弦所围的图形面积的最小值。

\begin{hintbox}
把过焦点的弦参数化，先写面积，再对参数求最小值。
\end{hintbox}
\end{exercisebox}

\begin{exercisebox}
\textbf{习题 3 \badge{弧长}}

求下列曲线的弧长：

\[
      \begin{aligned}
      &(1)\ y=x^{3/2},\quad0\le x\le4;\\
      &(2)\ x=\frac{y^2}{4}-\frac{\ln y}{2},\quad1\le y\le e;\\
      &(3)\ y=\ln\cos x,\quad0\le x\le a<\frac{\pi}{2};\\
      &(4)\ x=a\cos^3t,\ y=a\sin^3t,\quad0\le t\le2\pi;\\
      &(5)\ x=a(\cos t+t\sin t),\ y=a(\sin t-t\cos t),\quad0\le t\le2\pi;\\
      &(6)\ r=a(1-\cos\theta),\quad0\le\theta\le2\pi;\\
      &(7)\ r=a\theta,\quad0\le\theta\le2\pi;\\
      &(8)\ r=a\sin^3\frac{\theta}{3},\quad0\le\theta\le3\pi.
      \end{aligned}
      \]

\begin{hintbox}
识别曲线形式后直接选弧长公式；参数曲线先算 \(\sqrt{[x']^2+[y']^2}\)。
\end{hintbox}
\end{exercisebox}

\begin{exercisebox}
\textbf{习题 4 \badge{旋轮线弧长}}

在旋轮线的第一拱上，求分该拱的长度为 \(1:3\) 的点的坐标。

\begin{hintbox}
先写旋轮线参数方程的弧长函数 \(s(t)\)，再令 \(s(t)\) 等于总长的相应比例。
\end{hintbox}
\end{exercisebox}

\begin{exercisebox}
\textbf{习题 5 \badge{截面积体积}}

求下列几何体的体积：

\[
      \begin{aligned}
      &(1)\ \text{正椭圆台：上底长半轴 }a\text{、短半轴 }b,\text{ 下底长半轴 }A\text{、短半轴 }B,\text{ 高 }h;\\
      &(2)\ \text{椭球体 }\frac{x^2}{a^2}+\frac{y^2}{b^2}+\frac{z^2}{c^2}=1;\\
      &(3)\ \text{直圆柱面 }x^2+y^2=a^2\text{ 和 }x^2+z^2=a^2\text{ 所围几何体};\\
      &(4)\ \text{球面 }x^2+y^2+z^2=a^2\text{ 和直圆柱面 }x^2+y^2=ax\text{ 所围几何体}.
      \end{aligned}
      \]

\begin{hintbox}
统一使用截面积法。第 (3) 是经典 Steinmetz 体；第 (4) 可用柱坐标或按圆柱投影区域积分。
\end{hintbox}
\end{exercisebox}

\begin{exercisebox}
\textbf{习题 6 \badge{旋转体公式证明}}

证明以下旋转体的体积公式：

\[
      (1)\quad
      V=2\pi\int_a^b xf(x)\,dx,
      \]

其中 \(f(x)\ge0\) 连续，区域为 \(0\le a\le x\le b,\ 0\le y\le f(x)\)，绕 \(y\) 轴旋转一周。

\[
      (2)\quad
      V=\frac{2\pi}{3}\int_\alpha^\beta r^3(\theta)\sin\theta\,d\theta,
      \]

其中极坐标区域为 \(0\le\alpha\le\theta\le\beta\le\pi,\ 0\le r\le r(\theta)\)，绕极轴旋转一周。

\begin{hintbox}
第 (1) 用圆柱壳；第 (2) 把小扇形绕极轴旋转得到近似圆锥壳。
\end{hintbox}
\end{exercisebox}

\begin{exercisebox}
\textbf{习题 7 \badge{旋转体体积}}

求下列曲线绕指定轴旋转一周所围成的旋转体体积：

\[
      \begin{aligned}
      &(1)\ \frac{x^2}{a^2}+\frac{y^2}{b^2}=1,\text{ 绕 }x\text{ 轴};\\
      &(2)\ y=\sin x,\ y=0,\ 0\le x\le\pi,\quad (i)\text{ 绕 }x\text{ 轴},\ (ii)\text{ 绕 }y\text{ 轴};\\
      &(3)\ x=a\cos^3t,\ y=a\sin^3t,\ 0\le t\le\pi,\text{ 绕 }x\text{ 轴};\\
      &(4)\ x=a(t-\sin t),\ y=a(1-\cos t),\ 0\le t\le2\pi,\ y=0,\\
      &\qquad (i)\text{ 绕 }y\text{ 轴},\ (ii)\text{ 绕直线 }y=2a;\\
      &(5)\ x^2+(y-b)^2=a^2,\quad0<a\le b,\text{ 绕 }x\text{ 轴};\\
      &(6)\ r=a(1-\cos\theta),\text{ 绕极轴};\\
      &(7)\ r=ae^\theta,\quad0\le\theta\le\pi,\text{ 绕极轴};\\
      &(8)\ (x^2+y^2)^2=a^2(x^2-y^2),\text{ 绕 }x\text{ 轴}.
      \end{aligned}
      \]

\begin{hintbox}
绕 \(x\) 轴常用圆盘/垫片；绕 \(y\) 轴常用圆柱壳；极坐标绕极轴用习题 6(2) 的公式。
\end{hintbox}
\end{exercisebox}

\begin{exercisebox}
\textbf{习题 8 \badge{体积相等}}

将抛物线 \(y=x(x-a)\) 与 \(y=0\) 所界的区域在 \(x\in[0,a]\) 和 \(x\in[a,c]\) 的部分分别绕 \(x\) 轴旋转一周后，所得旋转体的体积相等，求 \(c\) 与 \(a\) 的关系。

\begin{hintbox}
分别写 \(\pi\int y^2dx\)，注意 \(y^2=x^2(x-a)^2\)。
\end{hintbox}
\end{exercisebox}

\begin{exercisebox}
\textbf{习题 9 \badge{半体积}}

记 \(V(\xi)\) 是曲线

\[
      y=\frac{\sqrt{x}}{1+x^2}
      \]

与 \(y=0\) 所界的区域在 \(x\in[0,\xi]\) 的部分绕 \(x\) 轴旋转一周所得旋转体体积，求常数 \(a\) 使得

\[
      V(a)=\frac12\lim_{\xi\to+\infty}V(\xi).
      \]

\begin{hintbox}
先写 \(V(\xi)=\pi\int_0^\xi \frac{x}{(1+x^2)^2}dx\)。
\end{hintbox}
\end{exercisebox}

\begin{exercisebox}
\textbf{习题 10 \badge{穿孔椭球}}

将椭圆

\[
      \frac{x^2}{a^2}+\frac{y^2}{b^2}=1
      \]

绕 \(x\) 轴旋转一周围成一个旋转椭球体，再沿 \(x\) 轴方向用半径为 \(r\ (r<b)\) 的钻头打一个穿心的圆孔，剩下的体积恰为原来椭球体体积的一半，求 \(r\) 的值。

\begin{hintbox}
剩余体积可看成在 \(|y|\ge r\) 的部分绕轴留下的体积；先求孔的 \(x\) 方向范围。
\end{hintbox}
\end{exercisebox}

\begin{exercisebox}
\textbf{习题 11 \badge{面积最值与体积}}

设直线 \(y=ax\ (0<a<1)\) 与抛物线 \(y=x^2\) 所围成图形面积为 \(S_1\)，且它们与直线 \(x=1\) 所围成图形面积为 \(S_2\)。

(1) 试确定 \(a\) 的值，使得 \(S_1+S_2\) 达到最小，并求出最小值；

(2) 求该最小值所对应的平面图形绕 \(x\) 轴旋转一周所得旋转体的体积。

\begin{hintbox}
交点是 \(x=0,a\)。先分区间写 \(S_1,S_2\)，再对 \(a\) 求最小。
\end{hintbox}
\end{exercisebox}

\begin{exercisebox}
\textbf{习题 12 \badge{微分方程与体积最值}}

设函数 \(f(x)\) 在闭区间 \([0,1]\) 上连续，在开区间 \((0,1)\) 上大于零，并满足

\[
      xf'(x)=f(x)+\frac{3a}{2}x^2\qquad(a\text{ 为常数}).
      \]

进一步，假设曲线 \(y=f(x)\) 与直线 \(x=1\) 和 \(y=0\) 所围的图形 \(S\) 的面积为 \(2\)。

(1) 求函数 \(f(x)\)；

(2) 当 \(a\) 为何值时，图形 \(S\) 绕 \(x\) 轴旋转一周所得旋转体体积最小？

\begin{hintbox}
先把方程改成 \(\left(f(x)/x\right)'\) 的形式，再用面积条件确定常数。
\end{hintbox}
\end{exercisebox}

\begin{exercisebox}
\textbf{习题 13 \badge{旋转曲面面积}}

求下列旋转曲面的面积：

\[
      \begin{aligned}
      &(1)\ y^2=2px,\quad0\le x\le a,\text{ 绕 }x\text{ 轴};\\
      &(2)\ y=\sin x,\quad0\le x\le\pi,\text{ 绕 }x\text{ 轴};\\
      &(3)\ \frac{x^2}{a^2}+\frac{y^2}{b^2}=1,\text{ 绕 }x\text{ 轴};\\
      &(4)\ x=a\cos^3t,\ y=a\sin^3t,\quad0\le t\le\pi,\text{ 绕 }x\text{ 轴};\\
      &(5)\ r=a(1-\cos\theta),\text{ 绕极轴};\\
      &(6)\ r^2=a^2\cos2\theta,\quad (i)\text{ 绕极轴},\ (ii)\text{ 绕射线 }\theta=\pi/2.
      \end{aligned}
      \]

\begin{hintbox}
统一记住：曲面面积微元 \(dS=2\pi\cdot\text{旋转半径}\cdot ds\)。
\end{hintbox}
\end{exercisebox}

\begin{exercisebox}
\textbf{习题 14 \badge{切线与曲面面积}}

设曲线 \(y=\sqrt{x-1}\)，过原点作其切线，求由该曲线、所作切线及 \(x\) 轴所围成的平面图形绕 \(x\) 轴旋转一周所得旋转体的表面积。

\begin{hintbox}
先求过原点切线的切点，再分别计算曲线段与直线段旋转形成的曲面面积。
\end{hintbox}
\end{exercisebox}

\begin{exercisebox}
\textbf{习题 15 \badge{柱面面积公式}}

证明：由空间曲线

\[
      x=x(t),\quad y=y(t),\quad z=z(t),\qquad t\in[T_1,T_2]
      \]

垂直投影到 \(Oxy\) 平面所形成的柱面的面积公式为

\[
      S=\int_{T_1}^{T_2}z(t)\sqrt{[x'(t)]^2+[y'(t)]^2}\,dt,
      \]

这里假设 \(x'(t),y'(t),z'(t)\) 在 \([T_1,T_2]\) 上连续，且 \(z(t)\ge0\)。

\begin{hintbox}
底边微元是投影曲线弧长 \(ds\)，柱面小片面积近似为高度 \(z(t)\) 乘以 \(ds\)。
\end{hintbox}
\end{exercisebox}

\begin{exercisebox}
\textbf{习题 16 \badge{指定点曲率}}

求下列曲线在指定点的曲率和曲率半径：

\[
      (1)\ xy=4,\text{ 在点 }(2,2);
      \qquad
      (2)\ x=a(t-\sin t),\ y=a(1-\cos t),\text{ 在 }t=\pi/2\text{ 对应的点}.
      \]

\begin{hintbox}
第 (1) 可隐函数求 \(y'\)、\(y''\)；第 (2) 用参数曲率公式。
\end{hintbox}
\end{exercisebox}

\begin{exercisebox}
\textbf{习题 17 \badge{曲率公式应用}}

求下列曲线的曲率和曲率半径：

\[
      \begin{aligned}
      &(1)\ y^2=2px\quad(p>0);\\
      &(2)\ \frac{x^2}{a^2}-\frac{y^2}{b^2}=1;\\
      &(3)\ x^{2/3}+y^{2/3}=a^{2/3}\quad(a>0);\\
      &(4)\ x=a(\cos t+t\sin t),\ y=a(\sin t-t\cos t)\quad(a>0).
      \end{aligned}
      \]

\begin{hintbox}
能参数化的优先参数化，避免隐函数二阶求导太重。
\end{hintbox}
\end{exercisebox}

\begin{exercisebox}
\textbf{习题 18 \badge{曲率圆}}

求曲线 \(y=\ln x\) 在点 \((1,0)\) 处的曲率圆方程。

\begin{hintbox}
先求曲率半径，再找主法线方向上的圆心。
\end{hintbox}
\end{exercisebox}

\begin{exercisebox}
\textbf{习题 19 \badge{极坐标曲率证明}}

设曲线的极坐标方程为 \(r=r(\theta),\theta\in[\alpha,\beta]\subset[0,2\pi]\)，且 \(r(\theta)\) 二阶可导。证明它在点 \((r,\theta)\) 处的曲率为

\[
      K=\frac{|r^2+2(r')^2-rr''|}{(r^2+(r')^2)^{3/2}}.
      \]

\begin{hintbox}
把 \(x=r\cos\theta,\ y=r\sin\theta\) 看成参数方程，再代入参数曲率公式。
\end{hintbox}
\end{exercisebox}

来源：教材第七章 §4（教材页 267-283）整理；公式用 LaTeX 渲染，课后题按教材顺序录入。
